\section{4.5.Correlator}\label{correlator}

\subsection{4.5. Módulo de
correlación}\label{muxf3dulo-de-correlaciuxf3n}

Una vez normalizados los resultados obtenidos por las herramientas de
análisis, EVMAudit aplica un proceso de correlación cuyo objetivo
consiste en identificar vulnerabilidades equivalentes detectadas por
distintas herramientas y agruparlas en un único hallazgo.

Este módulo constituye el núcleo de la contribución desarrollada en el
presente Trabajo Fin de Máster, ya que permite reducir la redundancia de
información y proporcionar una visión más estructurada de las
vulnerabilidades detectadas.

La necesidad de incorporar esta fase surge de una de las principales
limitaciones observadas durante el estudio del estado del arte. Aunque
herramientas como Slither y Mythril son capaces de detectar un gran
número de vulnerabilidades, cada una de ellas presenta los resultados
utilizando nomenclaturas y estructuras diferentes, generando además
múltiples duplicidades que dificultan la revisión manual por parte del
auditor.

Por este motivo, EVMAudit incorpora un mecanismo propio de correlación
encargado de consolidar la información procedente de ambas herramientas.

\subsubsection{4.5.1. Objetivos del proceso de
correlación}\label{objetivos-del-proceso-de-correlaciuxf3n}

Los objetivos perseguidos por este módulo son los siguientes:

\begin{itemize}
\tightlist
\item
  reducir la redundancia de información
\item
  agrupar vulnerabilidades equivalentes
\item
  incrementar la confianza en los hallazgos detectados
\item
  facilitar la interpretación de resultados
\item
  proporcionar una representación unificada de las vulnerabilidades.
\end{itemize}

La correlación no pretende eliminar completamente los falsos positivos
generados por las herramientas utilizadas, sino proporcionar una capa
adicional de análisis que facilite el trabajo posterior del auditor.

\subsubsection{4.5.2. Estrategia de
correlación}\label{estrategia-de-correlaciuxf3n}

El proceso implementado en EVMAudit se basa en una estrategia de
correlación mediante reglas.

Dos vulnerabilidades se consideran equivalentes cuando cumplen
simultáneamente las siguientes condiciones:

\begin{itemize}
\tightlist
\item
  afectan al mismo contrato
\item
  afectan a la misma función
\item
  poseen el mismo identificador SWC.
\end{itemize}

De este modo, la clave utilizada para la correlación puede expresarse de
la siguiente forma:

\begin{Shaded}
\begin{Highlighting}[]
\NormalTok{Contrato + Función + SWC}
\end{Highlighting}
\end{Shaded}

Esta aproximación permite agrupar vulnerabilidades que, aunque procedan
de herramientas diferentes, representan realmente el mismo problema de
seguridad.

La elección del identificador SWC como criterio común permite disponer
de una referencia independiente de la nomenclatura utilizada por cada
herramienta.

!TODO: insertar referencia cruzada a la Sección 2.X (SWC Registry).

!TODO: insertar pseudocódigo del algoritmo de correlación.

\subsubsection{4.5.3. Hallazgos confirmados y hallazgos
detectados}\label{hallazgos-confirmados-y-hallazgos-detectados}

Una vez realizado el proceso de agrupación, EVMAudit distingue entre dos
tipos de vulnerabilidades.

\paragraph{Hallazgos confirmados}\label{hallazgos-confirmados}

Se consideran confirmados aquellos hallazgos que han sido detectados por
más de una herramienta.

La coincidencia entre distintas técnicas de análisis aumenta la
confianza asociada a la vulnerabilidad y proporciona una mayor evidencia
de su existencia.

\paragraph{Hallazgos detectados}\label{hallazgos-detectados}

Corresponden a vulnerabilidades identificadas por una única herramienta.

Aunque presentan un menor nivel de confianza, siguen siendo incluidas en
el informe final con el objetivo de evitar la pérdida de información
potencialmente relevante.

Esta clasificación permite priorizar posteriormente la revisión manual
realizada por el auditor.

\subsubsection{4.5.4. Nivel de confianza}\label{nivel-de-confianza}

Con el objetivo de proporcionar información adicional sobre la
fiabilidad de cada hallazgo, EVMAudit asigna un nivel de confianza
basado en las herramientas que han detectado la vulnerabilidad.

La confianza asociada no representa una medida probabilística, sino un
indicador orientativo derivado del número de herramientas que coinciden
en el mismo hallazgo.

Este mecanismo permite distinguir aquellas vulnerabilidades respaldadas
por varias herramientas de aquellas detectadas únicamente por una de
ellas.

No obstante, la interpretación final continúa dependiendo del criterio
del auditor, ya que la correlación implementada no sustituye la
validación manual.

!TODO: revisar nomenclatura definitiva utilizada por el código
(\texttt{confidence\_score}).

\subsubsection{4.5.5. Gestión de
severidades}\label{gestiuxf3n-de-severidades}

Cuando una vulnerabilidad es detectada por varias herramientas, puede
ocurrir que cada una de ellas asigne niveles de severidad diferentes.

Para resolver esta situación, EVMAudit conserva la severidad más alta
reportada por las herramientas implicadas.

Este enfoque se adopta siguiendo un criterio conservador, priorizando la
revisión de aquellas vulnerabilidades que potencialmente puedan tener un
mayor impacto.

\subsubsection{4.5.6. Preservación de
evidencias}\label{preservaciuxf3n-de-evidencias}

Durante el proceso de correlación se mantiene la información procedente
de todas las herramientas que han participado en la detección.

De esta forma, cada hallazgo correlacionado conserva:

\begin{itemize}
\tightlist
\item
  las herramientas que lo han identificado
\item
  las evidencias originales asociadas
\item
  las localizaciones afectadas
\item
  la información contextual proporcionada por cada herramienta.
\end{itemize}

Esta decisión garantiza la trazabilidad del análisis y facilita la
revisión manual posterior.

\subsubsection{4.5.7. Beneficios del proceso de
correlación}\label{beneficios-del-proceso-de-correlaciuxf3n}

La incorporación de esta fase proporciona diversas ventajas:

\begin{itemize}
\tightlist
\item
  reducción de duplicidades
\item
  mejora de la legibilidad del informe final
\item
  aumento de la confianza en determinados hallazgos
\item
  simplificación del proceso de auditoría
\item
  independencia respecto a las herramientas utilizadas.
\end{itemize}

La correlación constituye, por tanto, uno de los principales elementos
diferenciadores de EVMAudit frente al uso individual de las herramientas
integradas.

!TODO: insertar diagrama del proceso de correlación.

!TODO: insertar ejemplo real antes y después de la correlación.

!TODO: insertar referencia a la Figura X.
